\textit{Идентификация варианта:}
\begin{itemize}
    \item Номер варианта: \underline{\hspace{4cm}}
    \item Схема нагружения: \underline{\hspace{4cm}}
    \item Сечение для определения прогиба: \underline{\hspace{4cm}}
\end{itemize}

\begin{longtblr}[
    caption = {Геометрические и физические параметры стержня},
    label = {tab:params},
]{
    colspec = { X[l,m] X[l,m] },
    hlines, vlines,
    rowhead = 1,
}
    \textit{Наименование параметра} & \textit{Значение (с указанием размерности)} \\
    Длина стержня $l$ & \\
    Ширина поперечного сечения $b$ & \\
    Высота поперечного сечения $h$ & \\
    Модуль нормальной упругости $E$ & \\
\end{longtblr}

\begin{longtblr}[
    caption = {Результаты измерений прогибов стержня при плоском поперечном изгибе},
    label = {tab:measurements},
]{
    colspec = { Q[c,m] X[c,m] X[c,m] X[c,m] },
    hlines, vlines,
    rowhead = 2,
}
    \SetCell[r=2]{c} \textit{Нагрузка $P$, Н} & \SetCell[c=3]{c} \textit{Показания индикатора, мкм} & & \\
    & \textit{Опыт 1} & \textit{Опыт 2} & \textit{Опыт 3} \\
    & & & \\
    & & & \\
    & & & \\
    & & & \\
    & & & \\
    & & & \\
\end{longtblr}